% Wave-16 S3/20 --- Borcherds Monster as a sibling specialisation
% Vol III, Calabi--Yau quantum groups (Raeez Lorgat)

\documentclass[11pt]{article}
\usepackage{amsmath,amssymb,amsthm}
\usepackage[margin=1in]{geometry}

\providecommand{\Z}{\mathbb{Z}}
\providecommand{\Q}{\mathbb{Q}}
\providecommand{\R}{\mathbb{R}}
\providecommand{\C}{\mathbb{C}}
\providecommand{\HH}{\mathbb{H}}
\providecommand{\Lat}{\Lambda}
\providecommand{\Leech}{\Lambda_{24}}
\providecommand{\IIlor}{\mathrm{II}_{25,1}}
\providecommand{\IIonel}{\mathrm{II}_{1,1}}
\providecommand{\gBKM}{\mathfrak{g}}
\providecommand{\mMon}{\mathfrak{m}}
\providecommand{\VLam}{V_\Leech}
\providecommand{\Vnat}{V^{\natural}}
\providecommand{\Aut}{\mathrm{Aut}}
\providecommand{\Mon}{\mathbb{M}}
\providecommand{\Co}{\mathrm{Co}}
\providecommand{\kBKM}{\kappa_{\mathrm{BKM}}}
\providecommand{\kch}{\kappa_{\mathrm{ch}}}
\providecommand{\kcat}{\kappa_{\mathrm{cat}}}
\providecommand{\Sp}{\mathrm{Sp}}
\providecommand{\Phii}{\Phi}
\providecommand{\Fcal}{\mathcal{F}}
\providecommand{\Hfa}{E_3\text{-hFA}}
\providecommand{\ClaimStatusTheorem}{[Theorem]}
\providecommand{\ClaimStatusConjectured}{[Conjectured]}

\theoremstyle{plain}
\newtheorem{theorem}{Theorem}
\newtheorem{proposition}[theorem]{Proposition}
\newtheorem{lemma}[theorem]{Lemma}
\theoremstyle{definition}
\newtheorem{construction}[theorem]{Construction}
\newtheorem{conjecture}[theorem]{Conjecture}

\title{Borcherds Monster as sibling specialisation of
$\mathcal{F}_{X_B}$: the $(\IIonel,C=\emptyset)$ cusp}
\author{Wave-16 S3/20 \\ Borcherds + Frenkel voices}
\date{}

\begin{document}
\maketitle

\paragraph{Primary denominator.} Borcherds 1992 proves
\[
  p^{-1}\!\!\prod_{m>0,\,n\in\Z}\!\!\bigl(1-p^m q^n\bigr)^{c(mn)}
  \;=\; j(p) - j(q),\qquad
  J(\tau)=j(\tau)-744=\sum_{n\ge-1} c(n)q^n,
\]
with $c(-1)=1,\;c(0)=0,\;c(1)=196\,884,\dots$ the McKay--Thompson
coefficients of $\Vnat$. The generalised Kac--Moody algebra
$\mMon$ whose simple roots are $\{(1,n): n\ge-1\}\subset\IIonel$
(multiplicity $c(n)$) has denominator $j(p)-j(q)$ (Borcherds 1992
Thm.~6.1). Its root lattice is the two-dimensional even Lorentzian
$\IIonel=U$, \emph{not} $\IIlor$; the rank-$26$ lattice belongs to
the Fake-Monster sibling $\gBKM_{\IIlor}$.

\paragraph{The CY$_3$ source.} \ClaimStatusConjectured{}
The Monster CY$_3$ is the Borcherds--Frenkel--Lepowsky--Meurman
orbifold
\[
  X_B \;=\; \bigl(T^{24}_\Leech \times E_\tau\bigr)/(\Z/2),
  \qquad \Z/2 \text{ acts as } (-1)_{T^{24}}\times(-1)_{E_\tau},
\]
with $T^{24}_\Leech=\R^{24}/\Leech$ the Leech torus and $E_\tau$
a generic elliptic curve. The fixed locus of $(-1,-1)$ is
$2^{24}\times 4$ points, each resolved by an $A_1$ crepant
resolution; the minimal resolution $X_B$ is a smooth compact
CY$_3$ with $\chi(\mathcal{O}_{X_B})=0$ (Kummer/FLM calculation,
FLM 1988 \S 12), hence $\kcat(X_B)=0$ (Künneth on products;
orbifold $E_2$-reduction does not alter $\chi(\mathcal{O})$).

\paragraph{$V^\natural$ as $E_3$-hFA shadow.} \ClaimStatusTheorem{}
The holomorphic VOA $\Vnat=(\VLam)^{\Z/2}\oplus(\VLam^T)^{\Z/2}$
of FLM 1988 (untwisted plus twisted sector) has central charge
$c=24$, character $J(\tau)$, and automorphism group $\Mon$ (FLM
Thm.~12.3.1; Griess 1982 for the algebra). Under the CY-to-chiral
functor $\Phii:\mathrm{CY}\text{-cat}_3\to\mathrm{ChirAlg}$, applied
to the derived category $D^b(X_B)$ equipped with its CY$_3$ Serre
structure, the target is $E_1$-chiral (Vol III dimensional
stratification, $d\ge3$). The $E_2$-factorisation structure sits
on the Drinfeld centre $Z(\mathrm{Rep}(\Phii(D^b X_B)))$, not on
$\Phii(D^b X_B)$ itself (the Vol~III $E_n$-hierarchy stratification). The
$E_3$-hFA \emph{shadow} --- the rank-$24$ holomorphic
factorisation algebra on $\R^2$ whose $E_3$-centre recovers the
full chiral structure --- is identified with $\Vnat$
by the Costello--Gwilliam dictionary: the holomorphic
factorisation of $E_\tau$-local operators in a topological
$(\Z/2)$-orbifold of a Leech-torus sigma model produces
$(\VLam)^{\Z/2}$ in the untwisted sector and the twisted-sector
module $(\VLam^T)^{\Z/2}$ by the Dijkgraaf--Witten orbifold
formula; their direct sum is $\Vnat$ by FLM.

\paragraph{Two-stage factorisation.} The universal family
$\Fcal_X$ on $\mathrm{Def}(X,\mathcal{L})$ sits over the period
domain of the transcendental lattice. For $X=X_B$ the relevant
period domain is
$\mathcal{D}_{\IIlor}$ (sibling of $\Delta_5$'s $\mathcal{D}_{3,2}$
and of $\Phi_{10}$'s $\mathcal{D}_{3,3}$). Two-stage factorisation
$\Phii = \Sp_{\Sigma_2,C}\circ\Fcal_X$:

\begin{construction}[Monster specialisation]
\label{cons:monster-specialisation}
Fix $X_B$ as above. The first stage
$\Fcal_{X_B}:\mathrm{Def}(X_B)\to\mathrm{VOA}^{\mathrm{fact}}$
assigns to each deformation the holomorphic factorisation algebra
$\Vnat_\Xi$ over the orbifold locus $\Xi$. The second stage
$\Sp_{\Sigma_2,C}$ specialises to
\[
  \Sigma_2 \;=\; \text{the Leech-fibre } T^{24}_\Leech,
  \qquad
  C \;=\; \emptyset \ (\text{point-cusp of } E_\tau),
\]
i.e.\ contracts $E_\tau$ to a nodal point and keeps only the
Leech-theta contribution. The resulting BKM-denominator identity is
\[
  \Sp_{T^{24}_\Leech,\,\mathrm{pt}}\!\bigl(\Fcal_{X_B}\bigr)
  \;=\; p^{-1}\prod_{m>0,n}\bigl(1-p^m q^n\bigr)^{c(mn)}
  \;=\; j(p)-j(q),
\]
where the twist variable $p$ records the Leech-orbifold twisted
sector and $q$ the $E_\tau$-modulus. The root lattice of the
resulting GKM is $\IIonel=U$ because the Leech part has collapsed
to the discrete McKay--Thompson grading $c(mn)$.
\end{construction}

\paragraph{Lattice siblings.} Monster and Igusa $\Delta_5$ siblings:
\begin{center}
\begin{tabular}{lll}
 & Monster $\mMon$ & Igusa $\gBKM_{\Delta_5}$ \\\hline
Root lattice & $\IIonel=U$ (sig.\ $(1,1)$)
            & $U\oplus U\oplus\langle-2\rangle$ (sig.\ $(3,2)$) \\
Envelope & $\gBKM_{\IIlor}$ (Fake Monster)
         & $\gBKM_{\Lat^{3,3}}$ (Pfaffian envelope, M3) \\
CY$_3$ & $(T^{24}_\Leech\times E)/\Z/2$ & $K3\times E$ \\
$(\Sigma_2,C)$ & (Leech fibre, $\mathrm{pt}$) & (K3 fibre, $E$) \\
Borcherds weight & $\kBKM(j(p){-}j(q))=0/2=0$ & $\kBKM(\Delta_5)=10/2=5$ \\
Frame shape & $j$-Hauptmodul & Gritsenko $\Delta_5$, $N{=}1$ \\
\end{tabular}
\end{center}
$j(p)-j(q)$ is modular weight $0$: a Borcherds product of weight
$c_F(0)/2$ with $c_F(0)=0$ (the input is $F=E_4^2 E_6/\Delta=j-744$,
whose $q^0$ Fourier coefficient in the Eichler--Zagier vector-valued
lift vanishes). This is the \emph{$N=0$ limit} of the
$N\in\{1,2,3,4,6\}$ Gritsenko series: the frame-shape index degenerates
to the Hauptmodul $j$ of level $1$, where $c_0(0)=0$, hence
$\kBKM=0$. The spectrum
\[
  \kBKM(\Mon,\,\Delta_5,\,\Phi_2,\,\Phi_3,\,\Phi_4,\,\Phi_6)
  \;=\;(0,\,5,\,2,\,1,\,1,\,1)
\]
extends the $K3\times E$ series $\{2,3,5,24\}$ by a sixth entry at
the $\IIonel$ cusp.

\paragraph{Envelope pattern (M3 refinement).}
The Wave-15 M3 envelope exhibits
$\Lat^{3,2}\hookrightarrow\Lat^{3,3}$ as codim-$1$ timelike
restriction, so that $\gBKM_{\Delta_5}$ sits as the $q^\perp$-fixed
subalgebra of the $\Lat^{3,3}$ GKM. The Monster sits analogously as
the $\IIonel\hookrightarrow\IIlor$ codim-$24$ restriction: project
$\IIlor=\Leech\oplus U\to U=\IIonel$ by quotienting by $\Leech$, and
the Fake-Monster GKM $\gBKM_{\IIlor}$ contains the Monster Lie
algebra $\mMon$ as its Leech-theta-averaged reduction
(Borcherds 1992 \S 8; Carnahan 2012 Thm.~1.4 for the generalised
moonshine automorphic-product closure under twisting).

\begin{theorem}[Leech-reduction]\label{thm:leech-reduction}
Let $\pi:\IIlor=\Leech\oplus U\twoheadrightarrow U=\IIonel$ be
projection onto the second summand. For each $(m,n)\in\IIonel$ with
$m>0$, $\sum_{\lambda\in\Leech}\mathrm{mult}_{\gBKM_{\IIlor}}(\lambda,m,n)
= c(mn)$, where $c$ are the McKay--Thompson coefficients. Hence
the pushforward $\pi_*\gBKM_{\IIlor}$ has root multiplicities matching
$\mMon$.
\end{theorem}
\begin{proof}
The Fake-Monster denominator (Borcherds 1990 Thm.~7.2) is
$\prod_{\alpha\in\IIlor^+}(1-e^\alpha)^{[1/\Delta(\tau)]_{-\alpha^2/2}}$.
Leech-summing at fixed $(m,n)\in U$ uses
$\theta_\Leech(\tau)/\Delta(\tau)=J(\tau)+24$ (Borcherds 1992
eq.~(8.2)), i.e.\ the Leech theta divided by the discriminant
reproduces the $j$-Hauptmodul up to the constant $24$ that matches
$c(0)=744-744=0$. The sum $\sum_\lambda$ at fixed $mn$-level is
$c(mn)$ by this identity.
\end{proof}

\paragraph{Sibling corollary.}
\begin{conjecture}[Many BKMs from one CY$_3$]
\label{conj:many-bkm-one-cy3}
There is a moduli $\mathcal{M}_{\mathrm{BKM}}$ whose generic point is
a CY$_3$ of the form $(\Sigma_2\times C)/G$ for $\Sigma_2$ a hyper-K\"ahler
surface ($K3$ or Leech-torus), $C$ a curve (elliptic or nodal), and
$G$ a finite group acting symplectically, such that the specialisation
$\Sp_{\Sigma_2,C}$ stratifies $\mathcal{M}_{\mathrm{BKM}}$ by
$(\Sigma_2,C)$-pairs with the following GKMs as images:
\[
  \begin{array}{l|l}
   (\Sigma_2,C) & \text{GKM} \\\hline
   (K3,\,E) & \gBKM_{\Delta_5}\;\;\kBKM=5\\
   (K3,\,\mathrm{nodal}) & \gBKM_{\Phi_N},\ N\in\{2,3,4,6\}\\
   (T^{24}_\Leech,\,E) & \gBKM_{\IIlor}\;(\text{Fake Monster})\\
   (T^{24}_\Leech,\,\mathrm{pt}) & \mMon\;(\text{Monster},\kBKM=0)\\
  \end{array}
\]
The Monster and $\gBKM_{\Delta_5}$ are sibling specialisations of
the same two-stage factorisation; the $\Lat^{3,2}\hookrightarrow
\Lat^{3,3}$ of M3 is the codim-$1$ Humbert restriction within the
$K3\times E$ stratum, and $\IIonel\hookrightarrow\IIlor$ is the
codim-$24$ Leech restriction within the $(T^{24}\times E)/\Z/2$
stratum.
\end{conjecture}

\paragraph{Status.} Borcherds 1992 Thm.~6.1 is theorem. FLM 1988
Thm.~12.3.1 ($V^\natural$ exists, $\Aut\Vnat=\Mon$) is theorem.
The identification $X_B=(T^{24}_\Leech\times E)/\Z/2$ as a smooth
CY$_3$ resolution producing $\Vnat$ as $E_3$-hFA shadow is
\ClaimStatusConjectured{}; a direct construction via
Costello--Gwilliam orbifold-factorisation for the Leech sigma model
would promote it to theorem. The two-stage factorisation
$(\Sp_{\Sigma_2,C}\circ\Fcal_X)$ is the Vol III operating conjecture
(Parts II, VI); Conjecture~\ref{conj:many-bkm-one-cy3} is its
specialisation to the Monster$/\Delta_5$ siblings.

\paragraph{Primary sources.} Borcherds 1986 \emph{Proc.\ Natl.\
Acad.\ Sci.}\ \textbf{83} (vertex algebras, Monster module);
Borcherds 1992 \emph{Invent.\ Math.}\ \textbf{109}
(monstrous moonshine, $j(p){-}j(q)$ denominator, Thm.~6.1);
Frenkel--Lepowsky--Meurman 1988 \emph{Vertex Operator Algebras and
the Monster} \S 12 ($\Vnat$ from Leech $\Z/2$-orbifold, Thm.~12.3.1);
Borcherds 1990 \emph{Invent.\ Math.}\ \textbf{109} (Fake Monster
denominator, $\IIlor$); Carnahan 2012 Thm.~1.4 (generalised
moonshine, automorphic closure under twisting); Conway--Norton
1979 \emph{Bull.\ LMS}\ \textbf{11} (monstrous moonshine conjectures).

\end{document}
